\begin{textsample}{NEG Dim 5 – global\_south\_2024\_05 – Score -1.00 – global\_south\_2024\_05/107911259.txt}  \label{ex:f5_neg_321}
Palestinians may be gratified to see American university campuses erupt in outrage over Israel ' s offensive in Gaza, but some in the embattled enclave are also wondering why no similar protests have hit the Arab countries they long viewed as allies.
Demonstrations have rocked US universities this week, with confrontations between students, counter-protesters and police, but while there have been some protests in Arab states, they have not been nearly as large or as vociferous.
Protesters are detained at the University of California Los Angeles, during a pro-Palestinian protest in Los Angeles, California, US, May 2. -- Reuters"We follow the protests on social media every day with admiration but also with sadness.
We are sad that those protests are not happening also in Arab and Muslim countries,"said Ahmed Rezik, 44, a father of five sheltering in Rafah in Gaza ' s south."Thank you students in solidarity with Gaza.
Your message has reached us.
Thank you students of Columbia.
Thank you students,"was scrawled across a tent in Rafah @ @ @ @ @ @ @ @ @ @ Israel ' s offensive.
A girl walks next to a tent sprayed with a message in solidarity with pro-Palestinian university students in Rafah in the southern Gaza Strip, May 2. -- Reuters Reasons for the comparative quiet on Arab campuses and streets may range from a fear of angering autocratic governments to political differences with Hamas and its Iranian backers or doubts that any protests could impact state policy.
American students at elite universities may face arrest or expulsion from their schools, but harsher consequences could await Arab citizens protesting without state authorisation.
Members of UCLA faculty stand on the frontlines as protesters stand together in the encampment after they were asked to leave by UCLA campus police. -- Reuters And US students may feel more motivated to protest as their own government backs and arms Israel, while even those Arab countries that have full diplomatic relations with it have been strongly critical of its military campaign.
When asked about the conflict, Arabs from Morocco to Iraq have consistently voiced fury at Israel ' s actions and @ @ @ @ @ @ @ @ @ @ Ramadan celebrations across the region last month.
Some rallies to support Palestinians have erupted, notably in Yemen where the Houthis have joined the conflict with strikes on shipping in the Red Sea.
And Arabs around the region have also shown their horror at the bombardment and support for their fellow Arabs in Gaza on social media, even if they have not taken to the streets.
But whatever the reason for the lack of public protests, some people in Gaza are now drawing unfavourable comparisons between the tumult in the United States and the public reaction they can see in other Arab countries."I ask Arab students to do what the Americans have done.
They should have done more for us than the Americans,"said Suha al-Kafarna, displaced by the bombardment from home in northern Gaza.
Public opinion In Egypt, which made peace with Israel in 1979 and where President Abdel Fattah al-Sisi has largely outlawed public protests, the authorities fear that demonstrations against Israel could later turn against the @ @ @ @ @ @ @ @ @ @ bombardment in October, some demonstrators veered off the agreed route and started chanting \textbf{anti-government} slogans, prompting arrests."One can not see the lack of large public protests against the war and the muted reaction on the Egyptian street in isolation from a broader context of crackdown on all forms of public protest and assembly,"said Hossam Bahgat, head of the Egyptian Initiative for Personal Rights.
At the American University in Cairo, security forces are less likely to intervene on campus and there have been some protests.
But a student activist there who requested anonymity said they could still face consequences for demonstrating."Being arrested here is nothing like being arrested in the US.
It ' s completely different,"he said, adding that there was"the factor of fear"preventing many from taking to the streets.
In Lebanon, where success in studies has become even more personally important to many young people after years of political and economic crises that have shrunk their shot at future prosperity, @ @ @ @ @ @ @ @ @ @ approached at campus protests in Beirut declined to be interviewed, saying they feared repercussions from university authorities.
The complex histories of Lebanon and other Arab states such as Jordan that host many Palestinian refugees also play into the question of public protests.
In Lebanon, some people blame Palestinians for triggering the 1975-90 civil war.
Others fear any overt displays of support for Palestinians might be hijacked by the Iran-backed Hezbollah, which has been trading fire with Israel since the start of the Gaza conflict."The Arab world is not reacting like Columbia or Brown ( US universities ) because they don't have the luxury to do so,"said Makram Rabah, a history professor at the American University of Beirut.
A sign adorns a tent at a protest encampment in support of Palestinians in Gaza. -- Reuters Besides, he added, with public opinion already largely backing the Palestinian cause it was not clear what protests there would achieve."The dynamics of power and the way you @ @ @ @ @ @ @ @ @ @ compared to the US,"he said.
Header image : Demonstrators react during a protest at an encampment in support of Palestinians at the University of California Los Angeles, in Los Angeles, California, US on May 2. -- Reuters

% matched lemmas: anti-government
\end{textsample}
